\subsection{The exact image for one fixed family of local coordinates}

\begin{statement}{A.132}{Proved}
{All-place finite matching gives a unique completed coordinate}
For each finite prime \(p\), let
\(b_p\in\widehat{\Q_p^\times}\), and choose a real sign
\(\epsilon\in\{1,-1\}=\widehat{\R^\times}\).
Write \(b_{p,N}\) for its image in
\(\Q_p^\times/(\Q_p^\times)^N\), and use the analogous
real image of \(\epsilon\).
The following actual diagonal maps are injective:
\[
\begin{aligned}
 \lambda_N:\Q^\times/(\Q^\times)^N
   &\longrightarrow\prod_{p<\infty}
                       \Q_p^\times/(\Q_p^\times)^N,\\
 \widehat\lambda:\widehat{\Q^\times}
   &\longrightarrow\prod_{p<\infty}\widehat{\Q_p^\times}.
\end{aligned}
\tag{A.132.1}
\]
There is an exact equality of images
\[
\begin{aligned}
 &\operatorname{im}\bigl(\widehat{\Q^\times}
       \longrightarrow\{\pm1\}\times
                          \prod_{p<\infty}\widehat{\Q_p^\times}\bigr)\\
 &\quad=\bigl\{(\epsilon,(b_p)):\ \forall N\ge1\ \exists a_N\in\Q^\times,\\
 &\hspace{30mm}[a_N]_p=b_{p,N}\ (p<\infty),\quad
         [a_N]_\infty=[\epsilon]_\infty\bigr\}.
\end{aligned}
\tag{A.132.2}
\]
On the right, the \emph{same fixed} \(b_p\) is used at every
exponent and every place.
Every finite class \([a_N]\) is unique, and their reductions
are automatically compatible.  Thus there is no further
inverse-limit obstruction after the all-place finite
matching in (A.132.2).
If no real sign is prescribed, project the equality to
the finite places; all-place finite matching then determines
the one possible global sign.

There is also the exact character description of the same
image.  At every \(N\) the local tuple must lie in the
restricted product
\[
 {\mathbb P}^1_N=
 \prod_v' H^1(\Q_v,\mu_N)
       \quad\text{with respect to }H^1_{\rm nr}(\Q_p,\mu_N)
                      \text{ at }p\nmid N,
\tag{A.132.3}
\]
and it must satisfy
\[
 \sum_v\operatorname{inv}_v(b_{v,N}\cup\chi_v)=0
                  \quad\text{for every }\chi\in H^1(\Q,\Z/N).
\tag{A.132.4}
\]
The real component in this formula is the chosen sign.
These two conditions, for all \(N\), are equivalent to
(A.132.2), not only necessary conditions.

For ordinary \(b_p\in\Q_p^\times\), put \(e_p=v_p(b_p)\in\Z\).
The restricted-product conditions are exactly
\[
       \text{for every }N,\quad N\mid e_p
                             \text{ for all but finitely many }p.
\tag{A.132.5}
\]
Under this condition define the actual convergent local products
\[
 U_q=\epsilon\,q^{e_q}\prod_{p\ne q}p^{e_p}
                         \quad\text{in }\Q_q^\times .
\tag{A.132.6}
\]
Then the exact image criterion is \(b_q=U_q\) for every
finite prime \(q\).  The only possible completed global
coordinate is the class with sign \(\epsilon\) and valuation
vector \((e_p)_p\) in A.10.
\end{statement}
\begin{proof}
If a rational \(a\) is an \(N\)th power in every \(\Q_p\),
then \(N\mid v_p(a)\) for every prime \(p\).
Unique factorization gives \(a=\delta r^N\), with
\(r\in\Q_{>0}\) and \(\delta\in\{1,-1\}\).
For odd \(N\) the sign is itself an \(N\)th power.
For even \(N\), a negative sign would make \(-1\)
a square in \(\Q_2\), which is impossible because an
odd square is \(1\) modulo eight.
Thus \(a\in(\Q^\times)^N\).
Applying this argument to the quotient of two rational
representatives proves injectivity of \(\lambda_N\).
Applying it at every finite modulus proves injectivity
of \(\widehat\lambda\).
This proof uses all finite places, including the place
at each prime appearing in \(a\); it does not require
the doubled-exponent density-one argument of A.76.

Suppose the right side of (A.132.2) holds.
For \(N\mid M\), the rational quotient \(a_M/a_N\)
is an \(N\)th power in every \(\Q_p\), because both
local classes are the prescribed \(b_{p,N}\).
The just-proved injection therefore gives
\[
                 [a_M]\bmod N=[a_N].
\tag{A.132.7}
\]
The classes define one element of
\(\varprojlim_N\Q^\times/(\Q^\times)^N\).
Its localization is the given \(b_p\) at every finite
level, and hence in the full local completion.
Its sign is \(\epsilon\).
Conversely an element of the global completion has
rational representatives at every finite quotient,
and the defining localization maps give exactly the
right side of (A.132.2).
Uniqueness follows from (A.132.1).
Without a prescribed real sign, the class at modulus
two has one sign, and (A.132.7) forces that sign at
every even modulus; the same construction applies.

For the character formulation, Kummer theory identifies
\(\Q^\times/(\Q^\times)^N\) with \(H^1(\Q,\mu_N)\)
and the local quotients with the corresponding local
cohomology groups.  Their Cartier dual is the literal
constant module \(\Z/N\).
The finite-coefficient Poitou--Tate exact sequence is
exact at the middle term of
\[
 H^1(\Q,\mu_N)\longrightarrow {\mathbb P}^1_N
       \longrightarrow H^1(\Q,\Z/N)^\vee,\qquad
 (b_v)\longmapsto
       \left(\chi\longmapsto
                   \sum_v\operatorname{inv}_v(b_v\cup\chi_v)\right).
\tag{A.132.8}
\]
Here \({}^\vee=\operatorname{Hom}_{\rm cont}(-,\Q/\Z)\);
at real places the usual Tate modification has the
same degree-one group.  This is the precise
restricted-product localization sequence used in
A.55 and in
\cite[I, Theorem 4.10(b) and Corollary 4.15]{BIR_MilneADT}.
The sum is finite: outside a finite set both classes
are unramified, and their cup product inflates from
the procyclic residue group, whose degree-two torsion
cohomology is zero.
Exactness of (A.132.8) proves the equivalence of
(A.132.3)--(A.132.4) with all-place matching at that
finite exponent.  The preceding compatibility proof
then proves the asserted completed equality.

At \(p\nmid N\), the unramified Kummer subgroup is
the image of \(\Z_p^\times\).
Indeed the exact valuation decomposition identifies
its quotient with \(\Z/N\), and taking an \(N\)th root
of a unit is etale over the unramified integer ring.
For an ordinary \(b_p\), being in this subgroup is
therefore exactly \(N\mid e_p\).
The finitely many primes dividing \(N\) do not affect
the restricted-product condition.  This proves
(A.132.5).

Fix \(q\) and a positive integer \(j\).
The finite group \((\Z/q^j)^\times\) has finite exponent,
say \(h_{q,j}\).
Condition (A.132.5) says that \(h_{q,j}\mid e_p\)
for all but finitely many \(p\).
For \(p\ne q\), every such factor satisfies
\(p^{e_p}\equiv1\pmod{q^j}\), also for negative exponents.
The product of these units therefore converges in
\(\Z_q^\times\), independently of the order of its
factors.  This proves the ordinary convergence in
(A.132.6).
At a finite power quotient, all but finitely many
global factors are powers of the modulus.
The localization of the A.10 class with sign
\(\epsilon\) and valuation vector \((e_p)\) is
therefore exactly the product (A.132.6).
Its valuation vector is forced by localization,
so no other global class can work.
Finally \(\Q_q^\times\to\widehat{\Q_q^\times}\)
is injective by its valuation and compact-unit
decomposition in A.10.
Equality in local completions is consequently the
literal equality \(b_q=U_q\), proving the last assertion.
\end{proof}

\paragraph{The actual finite restriction maps in the initialization problem.}
For one fixed family of ordinary old centres \(t_p\),
set \(b_p=t_p-c\).
At exponent \(N\), let \(\mathcal L_N\) be the set of
all-place ordinary old-point tuples that lift through
some augmented cover \(T-c=aY^N\), \(a\in\Q^\times\).
For a finite set of places \(S\), the map
\[
  \pi_S:\prod_v C_N(\Q_v)_{\{T\ne c\}}
                  \longrightarrow\prod_{v\in S}C_N(\Q_v)_{\{T\ne c\}}
\tag{A.132.9}
\]
is literal restriction of the local points.
A.85 determines when the prescribed finite tuple
belongs to \(\pi_S(\mathcal L_N)\).
All-place fixed-centre matching instead requires the
one entire tuple to belong to \(\mathcal L_N\).
The equalities (A.132.2) and (A.132.8) apply to this
second requirement, at all \(N\).
A.133 gives an explicit completed harmonic old tuple
whose finite restrictions belong to every such
\(\pi_S(\mathcal L_N)\), while the whole tuple fails
already the restricted-product condition at \(N=2\).
This identifies the exact map at which finite
controlled-point initialization ceases to imply
all-place matching.

The exact finite checks in
\path{code/fixed_local_escaping_support_certificate.py}
verify the valuation/sign injection, coefficient
transition identities, and actual local Euler-tail
congruences; its output is retained beside the code.
The inverse-limit equality and the Poitou--Tate
equivalence are established by the proof above.

